%% =====================================================================
%%  Unit 15 — Let's go out!  (Inviter et sortir)
%%  Pistes : 15.01–15.04 dialogues · 15.05–15.06 vocabulaire · 15.07–15.08 grammaire
%%           15.09 prononciation · 15.10–15.13 exercices (textes en annexe)
%% =====================================================================

\newcommand{\otIllusUnXV}{\otIllusDuo{movies/104/86/15/0.10,party/190/92/20/0.09,tonight/122/14/13/0.08,
  tickets/150/8/11/0.08,show/196/22/12/0.07,free/100/50/12/0.07}{sam}{Let's go out!}{lea}{I'd love to!}}

\unite{15}{Let's go out!}{Inviter et sortir}{otIllusUnXV}

\situation{Épisode 15}{Fin novembre. Sam a deux places pour un spectacle à Broadway : il propose à Léa
de l'accompagner. Le samedi suivant, dans l'escalier, Mrs.\ Lopez lui fait une invitation très
américaine : fêter \en{Thanksgiving} en famille.}

\begin{objectifs}
\item proposer une sortie avec \en{Let's…} ;
\item inviter poliment avec \en{Would you like to…?} ;
\item accepter avec enthousiasme, ou refuser poliment ;
\item employer les pronoms compléments : \en{me}, \en{him}, \en{her}, \en{us}, \en{them} ;
\item mettre la bonne intonation pour accepter ou refuser.
\end{objectifs}

\auprogramme
  {Proposer, inviter, accepter, refuser poliment}
  {\en{Let's…}, \en{Would you like to…?}\par\vspace{0.6mm}Les pronoms compléments}
  {Sorties, spectacles, rendez-vous : 25 expressions}
  {L'intonation pour accepter ou refuser}
  {Halloween et Thanksgiving}

\semaine{casque/otNavy/Dialogues/Écouter et lire les deux scènes,
  bulle/otGreenDark/Vocabulaire/Apprendre les 25 expressions,
  blocs/otViolet/Grammaire/Inviter et les pronoms compléments,
  sons/otOrangeDark/Sons et culture/L'intonation et les fêtes,
  crayon/otNavy/Exercices/12 exercices corrigés,
  coche/otGreenDark/Bilan/L'essentiel et le test express}

\clearpage
%% ---------------------------------------------------------------------
\rubrique{ecouter}{Sorties et invitations}

\begin{dialogue}{1}{Would you like to come?}{Mardi, 18 heures. Léa passe au café où travaille Sam,
  juste avant la fermeture.}{15.01}{15.02}
\dit{sam}{Léa! \hl{Are you free} on Friday evening?\note{1}}{Léa ! Tu es libre vendredi soir ?}
\dit{lea}{Yes, I am. Why?}{Oui. Pourquoi ?}
\dit{sam}{\hl{Would you like to see} a musical on Broadway? I have two tickets!\note{2}}{Tu veux voir une comédie musicale à Broadway ? J'ai deux places !}
\dit{lea}{Oh, \hl{I'd love to}! What's the show?}{Oh, avec plaisir ! C'est quel spectacle ?}
\dit{sam}{It's a new musical. My friend Ben is in it, and he gives \hl{me} free tickets.\note{3}}{C'est une nouvelle comédie musicale. Mon ami Ben joue dedans, et il me donne des places gratuites.}
\dit{lea}{\hl{That sounds great}! \hl{Let's have} dinner before the show.\note{4}}{Ça a l'air génial ! Allons dîner avant le spectacle.}
\dit{sam}{\hl{Great idea}! \hl{Let's meet} at six, in Times Square.}{Bonne idée ! Retrouvons-nous à six heures, à Times Square.}
\dit{lea}{Perfect. Text \hl{me} on Friday!}{Parfait. Envoie-moi un texto vendredi !}
\end{dialogue}

\begin{comprendre}
\pointnote{1} \en{\textbf{Are you free…?}} Avant d'inviter, on vérifie que la personne est disponible :
  \en{free} = « libre » (mais aussi « gratuit », comme dans \en{free tickets}).
\pointnote{2} \en{\textbf{Would you like to see…?}} La formule polie pour inviter : \en{Would you like}
  + \en{to} + base verbale. On répond \en{I'd love to!} (« avec plaisir ! »).
\pointnote{3} \en{\textbf{he gives me…}} \en{me} est un pronom complément : « me », « moi ». Après un
  verbe, \en{I} devient \en{me}.
\pointnote{4} \en{\textbf{Let's have dinner…}} \en{Let's} (= \en{let us}) + base verbale sert à
  proposer : « allons… », « si on… ? ».
\end{comprendre}

\begin{dialogue}{2}{Thanksgiving}{Samedi matin, dans l'escalier de l'immeuble. Mrs.\ Lopez revient du
  marché, les bras chargés.}{15.03}{15.04}
\dit{lopez}{Hi, Léa! What are you doing on Thursday?}{Bonjour, Léa ! Qu'est-ce que vous faites jeudi ?}
\dit{lea}{On Thursday? Nothing. Why?}{Jeudi ? Rien. Pourquoi ?}
\dit{lopez}{It's Thanksgiving! \hl{Would you like to have} dinner with \hl{us}?\note{5}}{C'est Thanksgiving ! Voulez-vous dîner avec nous ?}
\dit{lea}{Oh, that's very kind! \hl{I'd love to}. Can I help \hl{you}?}{Oh, c'est très gentil ! Avec plaisir. Je peux vous aider ?}
\dit{lopez}{Yes! Can you make a French dessert? My grandsons love \hl{them}!}{Oui ! Pouvez-vous faire un dessert français ? Mes petits-fils adorent ça !}
\dit{lea}{\hl{Sure}! I can make a chocolate cake.}{Bien sûr ! Je peux faire un gâteau au chocolat.}
\dit{lopez}{Perfect! Sofia, David and the kids come from Chicago every year. You can meet \hl{them}!\note{6}}{Parfait ! Sofia, David et les enfants viennent de Chicago tous les ans. Vous allez pouvoir faire leur connaissance !}
\dit{lea}{Great! What time is dinner?}{Super ! Le dîner est à quelle heure ?}
\dit{lopez}{It's at four. And bring your appetite!}{À quatre heures. Et venez avec de l'appétit !}
\end{dialogue}

\begin{comprendre}
\pointnote{5} \en{\textbf{with us}} Après une préposition (\en{with}, \en{for}, \en{to}), on emploie
  aussi le pronom complément : \en{we} devient \en{us}.
\pointnote{6} \en{\textbf{You can meet them}} \en{they} devient \en{them} : « les », « eux ». \en{meet}
  signifie ici « faire la connaissance de ».
\end{comprendre}

%% ---------------------------------------------------------------------
\rubrique{mots}{Les mots de l'unité}

Voici les 25 expressions de l'unité. Retenez surtout les réponses toutes faites : elles vous
serviront dans toutes les situations.

\begin{vocabulaire}{Les sorties}{15.05}
\categorie{Où aller ?}
\mot{go out}{ˌɡoʊ ˈaʊt}{sortir}
\mot{go to the movies}{ˌɡoʊ tə ðə ˈmuːviz}{aller au cinéma}
\mot{a movie}{ə ˈmuːvi}{un film}
\mot{a musical}{ə ˈmjuːzɪkəl}{une comédie musicale}
\mot{a show}{ə ˈʃoʊ}{un spectacle}
\mot{a theater}{ə ˈθiːətɚ}{un théâtre}
\mot{a party}{ə ˈpɑːrti}{une fête, une soirée}
\categorie{S'organiser}
\mot{a ticket}{ə ˈtɪkɪt}{un billet, une place}
\mot{free}{friː}{libre ; gratuit}
\mot{busy}{ˈbɪzi}{occupé}
\mot{meet}{miːt}{retrouver ; faire la connaissance de}
\mot{invite}{ɪnˈvaɪt}{inviter}
\mot{bring}{brɪŋ}{apporter, amener}
\mot{tonight}{təˈnaɪt}{ce soir}
\end{vocabulaire}

\begin{vocabulaire}{Inviter, accepter, refuser}{15.06}
\categorie{Proposer}
\mot{Let's go out!}{ˌlets ɡoʊ ˈaʊt}{Sortons !}
\mot{Would you like to come?}{wʊd juː ˌlaɪk tə ˈkʌm}{Tu veux venir ? Voulez-vous venir ?}
\mot{Are you free on Friday?}{ɑːr juː ˈfriː ɑːn ˈfraɪdeɪ}{Tu es libre vendredi ?}
\mot{Let's meet at six.}{ˌlets ˈmiːt ət ˈsɪks}{Retrouvons-nous à six heures.}
\categorie{Accepter}
\mot{I'd love to!}{aɪd ˈlʌv tuː}{Avec plaisir !}
\mot{Sure!}{ʃʊr}{Bien sûr ! Volontiers !}
\mot{Great idea!}{ˌɡreɪt aɪˈdiːə}{Bonne idée !}
\mot{That sounds great!}{ðæt ˌsaʊndz ˈɡreɪt}{Ça a l'air super !}
\categorie{Refuser poliment}
\mot{I'm sorry, I can't.}{aɪm ˈsɑːri aɪ ˈkænt}{Désolé, je ne peux pas.}
\mot{I'm busy.}{aɪm ˈbɪzi}{Je suis occupé.}
\mot{Maybe another time.}{ˌmeɪbi əˌnʌðɚ ˈtaɪm}{Une autre fois, peut-être.}
\end{vocabulaire}

%% ---------------------------------------------------------------------
\rubrique{grammaire}{Inviter et les pronoms compléments}

\otIntertitre[otViolet]{1 · Proposer, inviter, répondre}
\begin{center}\otTableau
\begin{tabular}{@{}l@{\hspace{8mm}}>{\enOn}l<{\enOff}@{\hspace{8mm}}l@{}}
\otTH{Pour\ldots} & \otTH{On dit\ldots} & \otTH{Français}\\\midrule
proposer (entre amis) & \textbf{Let's} go out! \textbf{Let's} have dinner! & sortons ! allons dîner !\\
inviter poliment      & \textbf{Would you like to} come?                  & tu veux / voulez-vous venir ?\\
accepter              & I'd love to! · Sure! · Great idea!                & avec plaisir ! volontiers !\\
refuser poliment      & I'm sorry, I can't. I'm busy.                     & désolé, je ne peux pas\\\bottomrule
\end{tabular}\end{center}

\begin{regle}[\en{Would you like to…?} ou \en{Do you like…?}]
\en{Do you like dancing?} demande si l'on aime danser \textbf{en général} (unité 7). \en{Would you like
to dance?} est une \textbf{invitation}, pour maintenant ou pour bientôt : « tu veux danser ? ». Après
\en{Let's} et \en{would you like to}, le verbe est toujours à la base verbale.
\end{regle}

\blocaudio{15.07}{Écoutez et répétez}{Quatre invitations et leurs réponses.}

\otIntertitre[otViolet]{2 · Les pronoms compléments}
\begin{center}\otTableau
\begin{tabular}{@{}>{\enOn}c<{\enOff}@{\hspace{10mm}}>{\enOn\otSemi\color{otViolet}}c<{\enOff}@{\hspace{10mm}}>{\enOn}l<{\enOff}@{\hspace{10mm}}l@{}}
\otTH{Sujet} & \otTH{Complément} & \otTH{Exemple} & \otTH{Français}\\\midrule
I    & me   & Call \textbf{me}!            & appelle-moi !\\
you  & you  & I can see \textbf{you}.      & je te vois\\
he   & him  & I like \textbf{him}.         & je l'aime bien\\
she  & her  & Come with \textbf{her}.      & viens avec elle\\
it   & it   & I love \textbf{it}!          & j'adore ça !\\
we   & us   & Have dinner with \textbf{us}! & dîne avec nous !\\
they & them & Invite \textbf{them}!        & invite-les !\\\bottomrule
\end{tabular}\end{center}

\begin{regle}[Où les placer ?]
Le pronom complément se place \textbf{après} le verbe (\en{I love her}), ou après une préposition
(\en{with me}, \en{for us}, \en{to them}). Jamais devant le verbe, comme en français : « je
\textbf{l'}aime » se dit \en{I love \textbf{her}}.
\end{regle}

\blocaudio{15.08}{Écoutez et répétez}{Six phrases avec un pronom complément.}

\begin{attention}[Les pièges des francophones]
\fauxjuste{Would you like come?}{Would you like to come?}
\fauxjuste{Let's to go out!}{Let's go out!}
\fauxjuste{I'd love!}{I'd love to!}
\fauxjuste{Call I tonight!}{Call me tonight!}
\fauxjuste{I her love.}{I love her.}
\fauxjuste{Come with we!}{Come with us!}
\end{attention}

%% ---------------------------------------------------------------------
\rubrique{sons}{L'intonation pour accepter ou refuser}

En anglais, la mélodie de la voix compte autant que les mots. Un \en{I'd love to!} prononcé à plat
semble poli… mais pas du tout enthousiaste ! Pour accepter, partez de haut et faites
\textbf{descendre} nettement la voix sur le mot important \otDescend. Pour refuser poliment, faites
d'abord \textbf{monter} la voix sur l'excuse \otMonte, puis redescendez. Retenez aussi : la voix monte
à la fin d'une question fermée (réponse \en{yes} ou \en{no}), elle descend après un mot interrogatif.

\begin{phrases}{15.09}{Écoutez et répétez en imitant la mélodie. Les flèches indiquent si la voix monte
\otMonte{} ou descend \otDescend.}
\phrase[Would you like to come?]{Would you like to come? \otMonte}{question fermée : la voix monte}
\phrase[What time?]{What time? \otDescend}{mot interrogatif : la voix descend}
\phrase[I'd love to!]{I'd \otDescend love to!}{enthousiasme : on part de haut}
\phrase[Great idea!]{Great i\otDescend dea!}{la voix plonge sur \textbf{dea}}
\phrase[I'm sorry, but I can't.]{I'm sorry \otMonte, but I can't. \otDescend}{refus poli : on monte, puis on descend}
\phrase[Maybe another time.]{Maybe another time. \otDescend}{voix douce et souriante}
\end{phrases}

%% ---------------------------------------------------------------------
\rubrique{culture}{Halloween et Thanksgiving}

\begin{culture}[Thanksgiving, la fête de famille]
Le quatrième jeudi de novembre, les Américains rentrent en famille pour \en{Thanksgiving}, en souvenir
du repas partagé en 1621 par les colons de Plymouth et les Amérindiens Wampanoag. Au menu : dinde
rôtie (\en{turkey}), farce (\en{stuffing}), sauce aux canneberges (\en{cranberry sauce}) et tarte au
potiron (\en{pumpkin pie}). À New York, la grande parade de Macy's et ses ballons géants ouvrent la
journée ; le lendemain, \en{Black Friday}, lance les achats de Noël.
\end{culture}

\begin{paire}
\begin{culture}[Halloween]
Le 31 octobre, les enfants se déguisent (\en{costumes}) et sonnent aux portes en criant
\en{Trick or treat!} (« des bonbons ou un sort ! »). On creuse des citrouilles éclairées, les
\en{jack-o'-lanterns}. À Manhattan, une immense parade costumée traverse Greenwich Village.
\end{culture}
\begin{culture}[Sortir à Broadway]
Broadway compte une quarantaine de grandes salles de spectacle. Astuce : le kiosque \en{TKTS} de
Times Square vend des places à prix réduit pour le jour même. Et après le spectacle, on attend les
acteurs à la sortie des artistes, la \en{stage door} !
\end{culture}
\end{paire}

%% ---------------------------------------------------------------------
\rubrique{pratique}{Exercices}

\exercice[ecrire,niveau=1]{Les sorties}{Écrivez l'expression anglaise.}
\begin{questions}[2][french]
\item une comédie musicale → \enb{\trou[34mm]{a musical}}
\item un billet, une place → \enb{\trou[34mm]{a ticket}}
\item une fête, une soirée → \enb{\trou[34mm]{a party}}
\item aller au cinéma → \enb{\trou[34mm]{go to the movies}}
\item occupé → \enb{\trou[34mm]{busy}}
\item libre → \enb{\trou[34mm]{free}}
\item ce soir → \enb{\trou[34mm]{tonight}}
\item sortir → \enb{\trou[34mm]{go out}}
\end{questions}

\exercice[ecrire,niveau=1]{Les pronoms compléments}{Remplacez les mots soulignés par un pronom
complément, puis récrivez la phrase.}
\begin{anglais}
\begin{questions}
\item Call \underLine{Sam} tonight! \hfill \trou[64mm]{Call him tonight!}
\item I can see \underLine{Emma}. \hfill \trou[64mm]{I can see her.}
\item Come with \underLine{Jack and me}! \hfill \trou[64mm]{Come with us!}
\item I don't like \underLine{the movie}. \hfill \trou[64mm]{I don't like it.}
\item Can you help \underLine{Léa and Sam}? \hfill \trou[64mm]{Can you help them?}
\item Mrs.\ Lopez invites \underLine{Léa}. \hfill \trou[64mm]{Mrs. Lopez invites her.}
\end{questions}
\end{anglais}

\exercice[lire,niveau=1]{Oui ou non ?}{Chaque réponse accepte-t-elle (A) ou refuse-t-elle (R)
l'invitation ?}
\begin{anglais}
\begin{questions}[2]
\item I'd love to! \hfill\trou[9mm]{A}
\item Sorry, I'm busy. \hfill\trou[9mm]{R}
\item Great idea! \hfill\trou[9mm]{A}
\item Maybe another time. \hfill\trou[9mm]{R}
\item Sure! What time? \hfill\trou[9mm]{A}
\item I'm sorry, I can't. \hfill\trou[9mm]{R}
\end{questions}
\end{anglais}

\exercice[ecrire,niveau=2]{Proposez !}{Écrivez une proposition avec les mots donnés.}
\begin{questions}
\item \en{Let's} / \en{have dinner} \hfill \enb{\trou[78mm]{Let's have dinner!}}
\item \en{Would you like} / \en{come to my party} \hfill \enb{\trou[78mm]{Would you like to come to my party?}}
\item \en{Let's} / \en{meet at seven} \hfill \enb{\trou[78mm]{Let's meet at seven.}}
\item \en{Would you like} / \en{see a musical} \hfill \enb{\trou[78mm]{Would you like to see a musical?}}
\item \en{Let's} / \en{go out tonight} \hfill \enb{\trou[78mm]{Let's go out tonight!}}
\end{questions}

\exercice[lire,niveau=2]{L'invitation du studio}{Emma a affiché cette invitation au studio. Lisez-la,
puis indiquez si chaque phrase est vraie (V) ou fausse (F).}
\begin{anglais}
\begin{tcolorbox}[enhanced,width=118mm,center,colback=otOrangeTint!45,colframe=otOrangeDark,boxrule=0.8pt,
  arc=2mm,left=6mm,right=6mm,top=4mm,bottom=4mm,halign=center,
  borderline={0.4pt}{1.4mm}{otOrangeDark!60},fontupper=\fontsize{10.6}{14}\selectfont]
{\otSemi\fontsize{8.4}{10}\selectfont\color{otOrangeDark}\otCaps[16]{You're invited!}}\par\vspace{1.6mm}
{\otBold\fontsize{22}{26}\selectfont\color{otInk}Friendsgiving at Bright Ideas}\par\vspace{1.4mm}
\otTuile[5mm]{etoile}{otOrange}\enspace\otTuile[5mm]{etoile}{otOrangeDark}\enspace\otTuile[5mm]{etoile}{otOrange}\par\vspace{2mm}
Thanksgiving is on Thursday, so let's celebrate together before the holiday!\par
\textbf{Wednesday, November 25th, 5 p.m. · in the studio kitchen}\par\vspace{1mm}
Please bring a dish from your country: let's try them all!\par
Music, games… and a big surprise for Jack.\par\vspace{1.4mm}
{\otSemi\fontsize{9}{11}\selectfont\color{otOrangeDark}RSVP to Emma by Monday}
\end{tcolorbox}
\end{anglais}
\par\vspace{1mm}\noindent{\fontsize{8.6}{10}\selectfont\color{otMuted}\textit{Friendsgiving} : un
Thanksgiving entre amis · \textit{celebrate} : fêter · \textit{a dish} : un plat · \textit{RSVP} :
répondez s'il vous plaît · \textit{by Monday} : avant lundi\par}
\begin{anglais}
\begin{questions}[2]
\item The party is on Thursday. \hfill\trou[9mm]{F}
\item It starts at five p.m. \hfill\trou[9mm]{V}
\item The party is in a restaurant. \hfill\trou[9mm]{F}
\item You need to bring food. \hfill\trou[9mm]{V}
\item You answer Emma by Monday. \hfill\trou[9mm]{V}
\item There's a surprise for Léa. \hfill\trou[9mm]{F}
\end{questions}
\end{anglais}

\exercice[lire,niveau=2]{Le bon ordre}{Remettez ce dialogue dans l'ordre : écrivez les lettres.}
\begin{anglais}
\begin{questions}[1][english]
\item[a.] Sure! What time?
\item[b.] Are you free on Saturday?
\item[c.] Yes, I am. Why?
\item[d.] Let's meet at eight, in front of the theater.
\item[e.] Would you like to go to the movies?
\item[f.] Great! See you on Saturday!
\end{questions}
\end{anglais}
\par\vspace{2mm}\noindent{\otSemi\fontsize{10}{12}\selectfont Ordre :}\hspace{4mm}%
\trou[9mm]{b}\hspace{3mm}\trou[9mm]{c}\hspace{3mm}\trou[9mm]{e}\hspace{3mm}\trou[9mm]{a}\hspace{3mm}\trou[9mm]{d}\hspace{3mm}\trou[9mm]{f}

\exercice[ecouter,niveau=1,piste=15.10]{Accepte ou refuse ?}{Écoutez les six petits dialogues. La
personne accepte-t-elle ou refuse-t-elle l'invitation ? Cochez.}
\begin{questions}[3]
\item \choix{accepte}\quad\choix{refuse}
\item \choix{accepte}\quad\choix{refuse}
\item \choix{accepte}\quad\choix{refuse}
\item \choix{accepte}\quad\choix{refuse}
\item \choix{accepte}\quad\choix{refuse}
\item \choix{accepte}\quad\choix{refuse}
\end{questions}

\exercice[ecouter,ecrire,niveau=2,piste=15.11]{Le message d'Emma}{Emma laisse un message sur la
messagerie de Léa. Écoutez et complétez.}
\begin{anglais}
\begin{questions}
\item Emma's plan is for \trou[30mm]{Saturday}.
\item At two o'clock, they go to the \trou[30mm]{museum}.
\item At seven, they have dinner in \trou[30mm]{Chinatown}.
\item Emma's cousin Harry is from \trou[30mm]{London}.
\item Léa needs to \trou[30mm]{call} Emma back.
\end{questions}
\end{anglais}

\exercice[ecouter,ecrire,niveau=2,piste=15.12]{Dictée à trous}{Sam invite Léa. Écoutez et complétez.}
\begin{anglais}
\noindent\begin{tabular}{@{}l@{\hspace{6mm}}l@{}}
\otNomDit{sam} & Are you \trou{free} on Friday?\\[1.8mm]
\otNomDit{lea} & Yes, I am. Why?\\[1.8mm]
\otNomDit{sam} & Would you \trou{like} to go to a concert?\\[1.8mm]
\otNomDit{lea} & I'd \trou{love} to! \trou{Where} is it?\\[1.8mm]
\otNomDit{sam} & In Brooklyn. \trou{Let's} \trou{meet} at the station at eight.\\[1.8mm]
\otNomDit{lea} & OK. Can you send \trou{me} the address?\\[1.8mm]
\otNomDit{sam} & Sure, I can text \trou{it} to you.\\
\end{tabular}
\end{anglais}

\exercice[ecrire,niveau=2]{Traduisez}{Traduisez en anglais.}
\begin{questions}[1][french]
\item Tu es libre samedi ? — Oui. \textit{(réponse courte)}\lignes{1}
\item Tu veux venir à ma fête ? — Avec plaisir !\lignes{1}
\item Allons au cinéma ce soir !\lignes{1}
\item Désolé, je ne peux pas. Je suis occupé.\lignes{1}
\item Retrouvons-nous à huit heures, devant le théâtre.\lignes{1}
\item Appelle-moi demain !\lignes{1}
\item J'aime bien Sam. Je l'invite.\lignes{1}
\item Viens avec nous !\lignes{1}
\end{questions}

\exercice[ecrire,niveau=3]{Une invitation}{Écrivez un message (cinq phrases) pour inviter un ami : dites
quoi, quel jour, à quelle heure et où vous retrouver.}
\lignes{5}

\exercice[ecouter,parler,niveau=3,piste=15.13]{À vous de parler !}{Jack vous invite. Répondez
\textbf{à voix haute} pendant chaque pause.}
\begin{center}\otTableau
\begin{tabular}{@{}>{\otSemi\fontsize{8.6}{10}\selectfont\color{otNavy}}l@{\hspace{3mm}}>{\enOn}p{70mm}<{\enOff}@{\hspace{6mm}}p{78mm}@{}}
 & \otTH{Jack dit\ldots} & \otTH{Vous\ldots}\\\midrule
1. & Are you free on Saturday? & répondez par une réponse courte.\\
2. & Would you like to see a hockey game? & acceptez avec enthousiasme.\\
3. & Let's meet at six. Is that OK? & acceptez et demandez où.\\
4. & Would you like to have dinner after the game? & refusez poliment.\\
5. & No problem! Can I call you tomorrow? & acceptez.\\\bottomrule
\end{tabular}\end{center}

%% ---------------------------------------------------------------------
\rubrique{bilan}{Faites le point}

\begin{essentiel}
\begin{minipage}[t]{0.54\linewidth}
\otEssTitre{Les phrases clés}
\renewcommand{\arraystretch}{1.32}
\begin{tabular}{@{}>{\enOn\otSemi\fontsize{9.4}{11}\selectfont}l<{\enOff}@{\hspace{3mm}}>{\fontsize{9.2}{11}\selectfont\color{white!78!otNavy}}l@{}}
Are you free on Friday?    & Tu es libre vendredi ?\\
Would you like to come?    & Tu veux venir ?\\
Let's go out!              & Sortons !\\
I'd love to!               & Avec plaisir !\\
I'm sorry, I can't.        & Désolé, je ne peux pas.\\
Call me!                   & Appelle-moi !\\
\end{tabular}
\end{minipage}\hfill
\begin{minipage}[t]{0.41\linewidth}
\otEssTitre{Inviter}
{\fontsize{9.2}{11.6}\selectfont \en{Let's} + base verbale · \en{Would you like to} + base verbale\par}
\otEssTitre{Les pronoms compléments}
{\fontsize{9.2}{11.6}\selectfont \en{me}, \en{you}, \en{him}, \en{her}, \en{it}, \en{us}, \en{them} :
après le verbe ou la préposition\par}
\otEssTitre{Prononciation}
{\fontsize{9.2}{11.6}\selectfont question fermée : la voix monte · enthousiasme : la voix part de
haut et descend\par}
\end{minipage}
\end{essentiel}

\begin{jesais}
\item proposer une sortie avec \en{Let's…} ;
\item inviter poliment avec \en{Would you like to…?} ;
\item accepter avec enthousiasme, ou refuser poliment ;
\item employer les pronoms compléments \en{me}, \en{him}, \en{her}, \en{us}, \en{them} ;
\item parler de sorties : cinéma, spectacle, fête ;
\item mettre la bonne intonation dans une invitation et dans une réponse.
\end{jesais}

\clearpage
\begin{testexpress}
\qcm[b]{« Tu veux venir ? » (invitation) se dit :}{\en{Do you like to come?}}{\en{Would you like to come?}}{\en{Would you like come?}}
\qcm[a]{« Sortons ! » se dit :}{\en{Let's go out!}}{\en{Let's to go out!}}{\en{Let's going out!}}
\qcm[c]{Pour accepter avec enthousiasme :}{\en{Maybe another time.}}{\en{I'm busy.}}{\en{I'd love to!}}
\qcm[b]{« Appelle-moi ! » se dit :}{\en{Call I!}}{\en{Call me!}}{\en{Me call!}}
\qcm[c]{« Je l'aime bien » (Sam) se dit :}{\en{I him like.}}{\en{I like he.}}{\en{I like him.}}
\qcm[a]{« Viens avec nous ! » se dit :}{\en{Come with us!}}{\en{Come with we!}}{\en{Come us with!}}
\qcm[b]{« Je suis occupé » se dit :}{\en{I'm free.}}{\en{I'm busy.}}{\en{I'm full.}}
\qcm[a]{Dans \en{Would you like to come?}, la voix :}{monte à la fin}{descend à la fin}{reste plate}
\qcm[c]{\en{Thanksgiving} a lieu :}{le 31 octobre}{le 4 juillet}{le quatrième jeudi de novembre}
\qcm[b]{« \en{Trick or treat!} », c'est :}{un plat de Thanksgiving}{ce que disent les enfants à Halloween}{une salle de Broadway}
\end{testexpress}

\noindent{\fontsize{9.6}{12}\selectfont\color{otMuted}Niveau visé (CECRL, niveau A1) : « peut faire,
accepter ou refuser une invitation simple ». Les corrigés de l'unité sont page~\pageref{cor-u15-1} ;
les textes des pistes d'exercices, page~\pageref{trans-15.10}.\par}
